\section{Introduction}\label{sec:introduction}

Which coordinates of a decision problem can be hidden without changing the optimal action, and how hard is it to certify that claim exactly? We expected one answer. We found three fundamentally different ones. We study this exact-certification question across static, stochastic, and sequential regimes. For a decision problem $\mathcal{D}=(A,S,U)$ with $S = X_1 \times \cdots \times X_n$, a coordinate set $I$ is sufficient when agreement on $I$ forces agreement on the optimal-action set:
\[
s_I = s'_I \implies \Opt(s) = \Opt(s').
\]
This is the basic problem of exact relevance certification: how much of the state must remain visible to preserve the decision boundary, and what it takes to certify that the rest is irrelevant. The question resembles subset selection, but exact certification has a different quantifier structure and, as the paper shows, a different complexity landscape.

This is a certification problem rather than a forward-evaluation problem. Evaluating $U(a,s)$ or $\Opt(s)$ on a fully specified state is one task; proving that a family of coordinates is irrelevant requires ruling out every counterexample pair of states. That universal structure drives the static regime. A first structural point is worth stating up front: \textsc{Minimum-Sufficient-Set} looks existential-universal on its face, but in the static regime it collapses because a set is sufficient exactly when it contains every relevant coordinate. So the minimum sufficient set is just the relevant-coordinate set. At the same time, the optimizer quotient is not an arbitrary definition: it is the canonical coarsest abstraction that preserves optimal-action distinctions. \textsc{Anchor-Sufficiency} does not collapse in the same way, because the anchor assignment is itself part of the existential choice and must then satisfy a universal fiberwise condition. At this static boundary the paper touches classical reduct-style attribute reduction, but its main object is broader: the same exact certification predicate is then transported to stochastic conditioning and sequential policy preservation.

This distinction also explains why exactness matters. When a system replaces exact relevance certification with a heuristic simplification, the omitted relevance does not disappear; it is shifted into residual uncertainty that must be handled elsewhere in the system. We refer to that externalized burden as the \emph{simplicity tax}. The results of this paper show that this tax is not merely rhetorical: exact simplification can be computationally expensive, but inexact simplification does not erase decision-relevant structure; it relocates it.

\paragraph{Informal interpretive guide.}
The formal objects studied here admit a compact intuitive reading. A relevant coordinate is a point of \emph{leverage}: changing it can change the decision. The optimizer quotient isolates the decision \emph{signal}: it is the coarsest abstraction that still preserves optimal-action distinctions. The static collapse expresses \emph{inevitability}: in the static regime, the optimal simplification is forced by the structure of relevance rather than discovered by search. The simplicity tax is \emph{liability}: omitted exact relevance does not disappear, but must be handled elsewhere in the system. At the task level, sufficiency is \emph{observation}, anchor sufficiency is \emph{discovery}, and sequential certification is \emph{exposure}: omitted information may appear harmless now while still creating future liability elsewhere in the system.

The paper's central claim is that this is not a collection of isolated hardness facts. Once the same exact-certification question is tracked across richer input models, the difficulty escalates in a systematic way. Static certification is governed by counterexample exclusion, stochastic certification by majority-style conditional comparisons, and sequential certification by policy preservation across evolving information states. Read as a complexity matrix, adding probability lifts exact certification from coNP to PP, and adding time lifts it again to PSPACE. The optimizer quotient packages the same issue structurally: it is the coarsest abstraction of the state space that preserves optimal-action distinctions.

Throughout, we work with finite coordinate-structured decision problems under explicit encoding regimes. Section~\ref{sec:encoding} fixes the computational model used by every complexity claim.

\subsection{Contributions}

The paper is best read as one theorem package with four linked takeaways: a structural static collapse, a regime-sensitive complexity matrix, a sharp phase transition between tractable and hard exact certification, and a trilemma for exact certifiers in the hard regime.

\begin{enumerate}
\item \textbf{The Static Collapse Theorem.} In the static regime, the apparent quantifier complexity of exact minimization simplifies sharply: sufficient sets are exactly the supersets of the relevant-coordinate set. This collapses \textsc{Minimum-Sufficient-Set} to a coNP problem rather than a genuine $\Sigma_2^P$ search problem, while \textsc{Anchor-Sufficiency} remains $\Sigma_2^P$-complete because its existential anchor choice does not collapse.

\item \textbf{Unified regime-sensitive complexity matrix.} Inside one common model of finite coordinate-structured decision problems, we formalize coordinate sufficiency, relevance, minimal sufficient sets, anchor sufficiency, and the optimizer quotient, and we show that the resulting exact certification problems organize into a single regime hierarchy: static \textsc{Sufficiency-Check} and \textsc{Minimum-Sufficient-Set} are coNP-complete, static \textsc{Anchor-Sufficiency} is $\Sigma_2^P$-complete, stochastic sufficiency is PP-complete with PP-hard anchor and minimum variants, and the sequential sufficiency, anchor, and minimum queries are PSPACE-complete.

\item \textbf{Sharp phase transition between tractable and hard exact certification.} We prove an explicit-state versus succinct dichotomy: logarithmic-size sufficient sets admit polynomial-time algorithms in the explicit regime, while linear-size sufficient sets in the succinct regime inherit ETH-conditioned exponential lower bounds. We also isolate six structured tractable subcases under standard restrictions.

\item \textbf{The trilemma of exact certification.} We transfer the hierarchy to exact simplification, isolate a witness-checking lower bound and approximation obstructions, and prove an integrity-competence impossibility theorem: no exact certifier can be simultaneously sound, complete, and polynomial-budgeted on the hard regime. This trilemma is a fundamental reliability constraint, and we show that omitted exact relevance is externalized rather than eliminated.

\item \textbf{Lean-certified reduction and finite-decision core.} The main reductions, size-bounded hardness packages, exact finite boolean deciders, counted finite-search procedures, and explicit-state step-counting wrappers for the regime-typed query predicates are mechanized in Lean 4. For the stochastic and sequential completeness claims, the artifact fully internalizes the reduction side together with the finite decision procedures; the standard PP/PSPACE membership arguments remain part of the paper text rather than end-to-end Turing-machine witness proofs in the current repository.
\end{enumerate}

\subsection{Quantifier Structure}

The complexity gap comes from the structure of the question. Insufficiency has short witnesses: two states that agree on the proposed coordinates but induce different optimal-action sets. Sufficiency, by contrast, asks for the absence of every such witness. So the object-level task is forward evaluation on a fully specified state, while the meta-level task is universal verification over counterexample pairs. This difference drives the coNP classification and explains why relevance identification is not just another forward-evaluation problem.

The static regime also contains the paper's main structural simplification. Although minimum sufficiency can be written with an outer existential quantifier over coordinate sets, the existential layer collapses once sufficiency is characterized as containment of all relevant coordinates. The stochastic and sequential regimes do not simplify in the same way: conditioning introduces majority-style comparisons over fibers, while temporal evolution introduces policy-level dependence on hidden information.

\subsection{Formalization Snapshot}

\begin{center}
\begin{tabular}{lr}
\hline
\textbf{Metric} & \textbf{Value} \\
\hline
Lines of Lean 4 backing cited content & \LeanReleaseLines \\
Theorems/lemmas in cited-content files & \LeanReleaseTheorems \\
Proof files backing cited content & \LeanReleaseFiles \\
\texttt{sorry} count in cited-content files & \LeanReleaseSorry \\
\hline
\end{tabular}
\end{center}

The artifact builds with \texttt{lake build}. It certifies the reduction-correctness lemmas, size-bounded hardness packages, exact finite deciders, counted-search witnesses, tractability statements, and explicit-state step-counting wrappers cited in the main text. The PP/PSPACE membership directions remain paper proofs rather than full Turing-machine formalizations.

\subsection{Paper Structure}

Section~\ref{sec:formal-setup} introduces the abstract setup and the encoding contract. Sections~\ref{sec:hardness}, \ref{sec:stochastic}, \ref{sec:sequential}, and \ref{sec:regime-hierarchy} develop the paper's core theorem package: the Static Collapse Theorem, the regime-sensitive complexity matrix, and the resulting hierarchy from static to stochastic to sequential certification. Section~\ref{sec:dichotomy} gives the encoding-sensitive dichotomy and the sharp phase transition for exact certification, and Section~\ref{sec:tractable} records tractable subcases. Section~\ref{sec:engineering-corollaries} then collects the main implications for exact simplification, the trilemma of exact certification, witness checking, approximation, and externalized relevance. Section~\ref{sec:related} situates the work within formalized complexity theory and decision-theoretic relevance.

\subsection{Artifact Availability}

The standalone Lean proof artifact and submission snapshot are archived at:
\begin{center}
\url{https://doi.org/10.5281/zenodo.18140966}
\end{center}
The proofs build with the Lean toolchain specified in \texttt{lean-toolchain}.
